% English translation of checked_pages/esq_119_044.tex; source: cote 119, site 44, PDF page 45, typescript page 42.
% AI-translated by GPT-6 Astra (OpenAI, ChatGPT Pro mode).
\scanpage{119/44}
\mspage{42}

corresponding to it must be very close, at least for $\unsure{\sourceglyph{119A}}$ that
can be generated by a finite ``system of generators'' S ---
but it is to be feared that the slightest interesting question one
might ask about this situation is currently out of
reach ($^1$).

These earlier reflections regained some relevance
for me with my later reflection on tame theories.
It seems to me, indeed, that it is possible to associate naturally
with a ``differentiable theory'' $\unsure{\sourceglyph{119A}}$ a tame theory $\hand{\mathcal{T}}$
(probably having the same field of constants), in such a way that
every $\unsure{\sourceglyph{119A}}$-manifold\frapparatus{V-FR-044-001} is automatically endowed with a $\hand{\mathcal{T}}$-tame structure,
and conversely that for every compact $\hand{\mathcal{T}}$-tame space $X$, one can
find a closed nowhere dense tame subset $Y$ in $X$, such that $X \smallsetminus Y$
comes from an $\unsure{\sourceglyph{119A}}$-manifold, and that moreover this $\unsure{\sourceglyph{119A}}$-manifold
structure is unique at least in the following sense: two such
structures coincide in the complement of a nowhere dense tame
subset $Y' \supset Y$ of $X$. The theory of dévissage of
stratified tame structures (discussed in the preceding paragraph),
in the case of smooth strata, should moreover raise
much more precise questions still about the comparison of tame
structures with structures of differentiable type (or rather,
$\mathbb{R}$-analytic type). I suspect that the kind of axiomatization proposed
here for the notion of a ``differentiable theory'' would provide a
natural setting for formulating such questions with all the precision
and generality one could wish for.


7. Since March of last year, thus for nearly
a year, the greater part of my energy has been devoted to
reflection on the \uline{foundations of noncommutative} (\uline{co})\uline{homological}
\uline{algebra}, or what ultimately amounts to the same thing, \uline{homotopical}
\uline{algebra}. These reflections have taken concrete form in a voluminous
bundle of typewritten notes, intended to form the first
volume (currently nearing completion) of a work in two

\margin{44}
